% =============================================================================
%  Puzzle Generator and Validator
% =============================================================================
\section{Puzzle Generator and Validator}
\label{sec:generator_validator}

\subsection{Overview}

The benchmark corpus is maintained by two complementary tools:

\begin{description}
  \item[\texttt{FutoshikiGenerator}] Creates valid, uniquely-solvable Futoshiki
    puzzles from scratch using a three-phase pipeline: backtracking fill,
    constraint decoration, and clue removal.
  \item[\texttt{validate\_file} / \texttt{validate\_path}] Verifies that a
    puzzle file (a) is syntactically well-formed, (b) has a valid initial
    state, (c) has \emph{exactly one} solution, and (d) if an expected file is
    provided, that the recorded solution matches the Z3-computed solution and
    shares the same constraint set.
\end{description}

Both tools use the Z3 SMT solver for uniqueness reasoning — the only place in
the project where an external constraint solver is used.

% =============================================================================
\subsection{File Format}
\label{sec:file_format}

All puzzle files share a single text format understood by both the generator
and the \texttt{Parser}:

\begin{verbatim}
N
r0c0,r0c1,...,r0c(N-1)          <- N grid rows
...
h_r0_0,...,h_r0_(N-2)           <- N horizontal-constraint rows (N-1 values)
...
v_r0_0,...,v_r0_(N-1)           <- N-1 vertical-constraint rows (N values)
\end{verbatim}

Encoding conventions:
\begin{itemize}
  \item Grid cells: \texttt{0} = empty, \texttt{1}--\texttt{N} = given clue value.
  \item Constraint cells: \texttt{0} = none, \texttt{1} = \texttt{<},
        \texttt{-1} = \texttt{>}.
  \item Lines beginning with \texttt{\#} are comments and are ignored by the
        parser.
\end{itemize}

The benchmark corpus stores paired files under
\texttt{resource/benchmark/input/} (puzzle with blanks) and
\texttt{resource/benchmark/expected/} (fully solved grid) using the same
filename so they can be matched by name.

% =============================================================================
\subsection{Parser}
\label{sec:parser}

\texttt{core/parser.py} converts a file into a \texttt{Puzzle} object.
It reads the file in a single pass, stripping blank lines and comments.

\begin{enumerate}
  \item \textbf{Line 0}: reads $N$; raises \texttt{ParseError} if absent or
        non-integer, or if $N < 2$.
  \item \textbf{Lines 1–$N$}: parses $N$ grid rows, each with $N$
        comma-separated integers in $\{0, \ldots, N\}$.
  \item \textbf{Lines $N+1$–$2N$}: parses $N$ horizontal-constraint rows,
        each with $N-1$ values in $\{-1, 0, 1\}$.
        Each non-zero entry produces an \texttt{InequalityConstraint}.
  \item \textbf{Lines $2N+1$–$3N-1$}: parses $N-1$ vertical-constraint rows,
        each with $N$ values in $\{-1, 0, 1\}$.
  \item \textbf{Post-parse validation}: checks all grid values are in
        $\{0, \ldots, N\}$.
\end{enumerate}

The total line count of a well-formed file is $1 + N + N + (N-1) = 3N$.
Any deviation raises a \texttt{ParseError} with the 1-based line number,
enabling precise error messages.

% =============================================================================
\subsection{Puzzle Generator}
\label{sec:generator}

\texttt{FutoshikiGenerator} produces a puzzle in three sequential phases.
The generator is seeded with a fixed integer so all 24 benchmark puzzles are
reproducible deterministically.

\subsubsection{Phase 1: Full Grid Construction (\texttt{generate\_full\_grid})}

A complete $N \times N$ Latin square is built by backtracking with
\emph{Minimum Remaining Values} (MRV) cell selection.

\paragraph{MRV cell selection (\texttt{find\_empty}).}
Instead of filling cells in row-major order, the generator always picks the
empty cell with the fewest legal values remaining:

\begin{enumerate}
  \item Iterate all empty cells; for each count the number of values in
        $\{1, \ldots, N\}$ that would not violate any row-uniqueness,
        column-uniqueness, or active inequality constraint.
  \item Two fast-path cuts:
        \begin{itemize}
          \item Options $= 0$: return immediately — this cell is a \emph{dead end};
                the caller will backtrack without wasting time on other cells.
          \item Options $= 1$: return immediately — a \emph{forced cell};
                no point searching further.
        \end{itemize}
  \item Otherwise track the cell with the minimum option count and return it
        after scanning all cells.
\end{enumerate}

\paragraph{Validity check (\texttt{is\_valid}).}
For a candidate value $v$ at position $(r,c)$, the check tests:
\begin{itemize}
  \item Row uniqueness: no other cell in row $r$ holds $v$.
  \item Column uniqueness: no other cell in column $c$ holds $v$.
  \item Left horizontal constraint: if $h[r][c-1] \ne 0$ and $\text{grid}[r][c-1] \ne 0$,
        the relation $\text{grid}[r][c-1] \;\square\; v$ must hold.
  \item Right horizontal constraint: if $h[r][c] \ne 0$ and $\text{grid}[r][c+1] \ne 0$,
        the relation $v \;\square\; \text{grid}[r][c+1]$ must hold.
  \item Upper/lower vertical constraints: analogous checks for $v[r-1][c]$ and
        $v[r][c]$.
\end{itemize}

Note that constraints are only enforced when the \emph{neighbour cell is already
filled}; unfilled neighbours are skipped to avoid over-constraining the search.

\paragraph{Fill loop.}
The fill is a standard recursive backtracker: pick an empty cell via MRV,
shuffle the value list (using the seeded \texttt{rng}), try each value in
random order, recurse, and undo on failure.
The final complete grid is saved as \texttt{solution\_grid}.

\subsubsection{Phase 2: Constraint Decoration (\texttt{add\_constraints})}

Each potential constraint position is included independently at random with
probability equal to the \emph{density} parameter.
The direction is derived directly from the solution:

\begin{align*}
  h[r][c] &= \begin{cases}
    1  & \text{if } \text{sol}[r][c] < \text{sol}[r][c+1] \\
    -1 & \text{otherwise}
  \end{cases}
  \quad \text{with probability } d \\
  v[r][c] &= \begin{cases}
    1  & \text{if } \text{sol}[r][c] < \text{sol}[r+1][c] \\
    -1 & \text{otherwise}
  \end{cases}
  \quad \text{with probability } d
\end{align*}

Since the direction is read from the known solution, every added constraint
is automatically satisfied by construction — no separate feasibility check
is needed at this phase.

The expected number of constraints added is:
\[
  \mathbb{E}[|\mathcal{I}|] = d \cdot \bigl(N(N-1) + N(N-1)\bigr) = 2dN(N-1)
\]

\begin{table}[h]
  \centering
  \caption{Constraint density parameters in the benchmark corpus.}
  \label{tab:density}
  \begin{tabular}{lccc}
    \toprule
    Difficulty & Density $d$ & $\mathbb{E}[|\mathcal{I}|]$ at $N=6$ & Fill ratio \\
    \midrule
    Easy   & 0.55 & 33.0 & 0.50 \\
    Medium & 0.45 & 27.0 & 0.35 \\
    Medium & 0.40 & 24.0 & 0.30 \\
    Hard   & 0.35 & 21.0 & 0.22 \\
    \bottomrule
  \end{tabular}
\end{table}

\subsubsection{Phase 3: Clue Removal (\texttt{create\_puzzle})}

Starting from the complete solution grid, the generator removes given-clue
cells one by one while preserving the \emph{unique-solution property}.

\paragraph{Stopping criterion.}
If a \texttt{target\_fill\_ratio} is specified, removal stops as soon as the
number of remaining filled cells reaches $\lceil N^2 \cdot r \rceil$ where
$r$ is the target ratio.
This avoids over-removing clues beyond the desired difficulty level.

\paragraph{Z3 uniqueness check.}
Before confirming each removal, the generator calls
\texttt{\_count\_solutions\_with\_base(solver, cells, grid, limit=2)}, which
asks the Z3 solver to find at most 2 solutions to the current partial grid:

\begin{enumerate}
  \item A single \texttt{Solver} and $N^2$ Z3 integer variables are built
        once at the start of Phase 3 (\texttt{\_build\_base\_solver}).
        The base solver encodes: domain bounds, row and column \texttt{Distinct}
        constraints, and all inequality constraints.
        This model is \emph{shared} across all uniqueness checks.

  \item For each candidate removal, a \texttt{solver.push()} saves the
        current state, and the remaining given-cell assignments are added as
        equality constraints.

  \item The solver is invoked twice using a \emph{solution-blocking} loop:
        after finding one solution, the negation of the full assignment
        ($\bigvee_{r,c} (\texttt{cell}_{r,c} \ne v_{r,c})$) is added to
        exclude it.
        If a second \texttt{sat} is returned, the puzzle is ambiguous.

  \item \texttt{solver.pop()} restores the base state for the next check.
\end{enumerate}

If the removal would leave $\ne 1$ solution, the cell is put back
(\texttt{self.grid[r][c] = temp}).
This guarantees that at termination \texttt{self.grid} is a valid puzzle
with exactly one solution.

\paragraph{Cell removal order.}
Cells are shuffled randomly (seeded) before the removal loop.
Row-major order would systematically expose more clues in later rows; random
order spreads the given cells more evenly.

\paragraph{Complexity.}
Each removal attempt costs one Z3 invocation.
In the worst case (no cells can be removed) there are $N^2$ attempts.
In practice most cells are removable, so the loop runs on the order of $N^2$
iterations with short Z3 queries.
The \texttt{solver.push()/pop()} design avoids rebuilding the base model on
every check, keeping each incremental query fast.

\subsubsection{Benchmark Corpus Specification}

The 24 benchmark puzzles are declared as a static list
\texttt{BENCHMARK\_SPECS} at module level.
Each entry records: \texttt{filename}, \texttt{size}, \texttt{seed},
\texttt{fill\_ratio}, and \texttt{density}.

\begin{table}[h]
  \centering
  \caption{Benchmark corpus structure.}
  \label{tab:corpus}
  \begin{tabular}{lcccc}
    \toprule
    Size & Easy & Medium (1) & Medium (2) & Hard \\
    \midrule
    $4 \times 4$ & \checkmark & \checkmark & \checkmark & \checkmark \\
    $5 \times 5$ & \checkmark & \checkmark & \checkmark & \checkmark \\
    $6 \times 6$ & \checkmark & \checkmark & \checkmark & \checkmark \\
    $7 \times 7$ & \checkmark & \checkmark & \checkmark & \checkmark \\
    $8 \times 8$ & \checkmark & \checkmark & \checkmark & \checkmark \\
    $9 \times 9$ & \checkmark & \checkmark & \checkmark & \checkmark \\
    \midrule
    Total & 6 & 6 & 6 & 6 \\
    \bottomrule
  \end{tabular}
\end{table}

Puzzles are indexed 01–24 in the order (size 4 easy, 4 medium1, 4 medium2,
4 hard, 5 easy, \ldots, 9 hard) and seeded with $1000 + \text{index}$
to ensure reproducibility across re-generation runs.

% =============================================================================
\subsection{Puzzle Repository}
\label{sec:repository}

\texttt{InMemoryPuzzleRepository} provides the application layer with two
operations:

\paragraph{Corpus listing (\texttt{list\_entries}).}
On first call, the repository scans \texttt{resource/benchmark/input/} for
\texttt{.txt} files matching the filename pattern
\texttt{puzzle\_NN\_NxN\_difficulty.txt}.
Each match produces a \texttt{PuzzleEntry} (frozen dataclass) recording the
file path, stem name, grid size, and difficulty string.
Entries are sorted by size ascending, then by difficulty order (easy $<$
medium $<$ hard), then by name.
The result is cached so subsequent calls pay no disk I/O.

\paragraph{On-demand generation (\texttt{generate}).}
The repository provides three named difficulty presets that map to
\texttt{(fill\_ratio, density)} pairs:

\begin{table}[h]
  \centering
  \caption{Difficulty presets for on-demand generation.}
  \label{tab:difficulty_presets}
  \begin{tabular}{lcc}
    \toprule
    Difficulty & Fill ratio & Density \\
    \midrule
    easy   & 0.55 & 0.50 \\
    medium & 0.35 & 0.40 \\
    hard   & 0.20 & 0.25 \\
    \bottomrule
  \end{tabular}
\end{table}

If no seed is supplied, a random integer in $[0, 999\,999]$ is chosen.
After running the three generator phases,
\texttt{\_generator\_to\_puzzle(gen)} converts the generator's internal
representation (plain Python lists of integers) to a \texttt{Puzzle}:

\begin{enumerate}
  \item Wrap \texttt{gen.grid} in a NumPy array (dtype \texttt{int}).
  \item Convert \texttt{gen.h\_const[r][c] \ne 0} entries to
        \texttt{InequalityConstraint} objects with \texttt{direction="<"}
        (when the value is $1$) or \texttt{">"} (when $-1$).
  \item Do the same for \texttt{gen.v\_const}.
  \item Construct and return a \texttt{Puzzle} whose post-init runs the
        constraint lookup map and cached cell list.
\end{enumerate}

% =============================================================================
\subsection{Validator}
\label{sec:validator}

\subsubsection{Validation Stages}

\texttt{validate\_file(path, expected\_dir)} runs four checks in sequence,
failing fast at the first error:

\begin{enumerate}
  \item \textbf{Parse.}
        \texttt{Parser().parse(path)} — raises \texttt{ParseError} on
        any syntactic problem (malformed line, wrong token count, out-of-range
        value).

  \item \textbf{Initial consistency (\texttt{\_initial\_validation}).}
        Checks the given cells in the puzzle (before any solving):
        \begin{itemize}
          \item No duplicate value in any row among given cells.
          \item No duplicate value in any column among given cells.
          \item No given cell pair that already violates a horizontal or vertical
                inequality constraint.
        \end{itemize}
        This is a fast $O(N^2 + |\mathcal{I}|)$ pass that catches obviously
        broken inputs before paying the cost of a Z3 invocation.

  \item \textbf{Uniqueness (\texttt{solve\_puzzle}).}
        Builds a Z3 model (\texttt{\_build\_z3\_model}) encoding:
        domain bounds, row/column \texttt{Distinct}, given-cell equalities,
        and all inequality constraints.
        Calls \texttt{solver.check()} up to \texttt{limit=2} times.
        \begin{itemize}
          \item 0 solutions → \texttt{"unsatisfiable puzzle"}.
          \item 1 solution → pass.
          \item 2 solutions → \texttt{"puzzle has multiple solutions"}.
        \end{itemize}

  \item \textbf{Expected-solution check} (optional, when \texttt{expected\_dir}
        is provided):
        \begin{itemize}
          \item The expected file must exist at
                \texttt{expected\_dir / path.name}.
          \item Its grid must equal the Z3-computed solution
                (\texttt{np.array\_equal}).
          \item Its constraint signature must match the input file's
                (\texttt{\_constraint\_signature} produces a sorted tuple of
                \texttt{(cell1, cell2, direction)} triples from all constraints).
        \end{itemize}
\end{enumerate}

\subsubsection{Z3 Model Construction (\texttt{\_build\_z3\_model})}

The model is built identically in both the validator and the generator (the
generator's \texttt{\_build\_base\_solver} is an equivalent reimplementation
that also stores the constraint matrices before the puzzle object exists).

\begin{lstlisting}[language=Python, caption={Z3 model core (simplified).}]
solver = Solver()
cells = [[Int(f"cell_{r}_{c}") for c in range(N)] for r in range(N)]

for r, c in product(range(N), range(N)):
    solver.add(cells[r][c] >= 1, cells[r][c] <= N)
for r in range(N):
    solver.add(Distinct(cells[r]))
for c in range(N):
    solver.add(Distinct([cells[r][c] for r in range(N)]))
for row, col, v in puzzle.get_given_cells():
    solver.add(cells[row][col] == v)
for constraint in puzzle.h_constraints + puzzle.v_constraints:
    r1, c1 = constraint.cell1
    r2, c2 = constraint.cell2
    if constraint.direction == "<":
        solver.add(cells[r1][c1] < cells[r2][c2])
    else:
        solver.add(cells[r1][c1] > cells[r2][c2])
\end{lstlisting}

The solution-counting loop adds a \emph{blocking clause} after each found
solution to exclude it from subsequent checks:
\[
  \bigvee_{r=0}^{N-1} \bigvee_{c=0}^{N-1}
  \left( \texttt{cell}_{r,c} \ne v_{r,c} \right)
\]
This is the standard \emph{all-solutions enumeration} technique for SMT
solvers.

\subsubsection{Batch Validation (\texttt{validate\_path})}

\texttt{validate\_path(path, expected\_dir)} accepts either a single file or
a directory.
For directories it uses \texttt{Path.rglob("*.txt")} to find all puzzle files
recursively, sorts them, and calls \texttt{validate\_file} on each.
Progress is printed with an ASCII progress bar.
The function returns a list of \texttt{ValidationResult} dataclasses and
exits with code 0 on full success or 1 on any failure.

\subsubsection{Design Choice: Z3 over Project Solvers}

The validator deliberately uses Z3 rather than one of the project's own
solvers for uniqueness checking.
The reasons are:

\begin{enumerate}
  \item \textbf{Independence}: a solver bug would not propagate into
        benchmark corruption; the two code paths are completely separate.
  \item \textbf{Completeness guarantee}: Z3's DPLL(T) engine is a complete
        decision procedure for the linear arithmetic theory used here.
        The project's constraint-propagation solvers may not terminate on
        all inputs or may classify unsatisfiable puzzles as failures rather
        than contradictions.
  \item \textbf{All-solutions counting}: enumerating up to 2 solutions with
        blocking clauses is natural in Z3 but would require significant
        bookkeeping in a backtracking solver.
\end{enumerate}

The trade-off is an external dependency and start-up overhead per call, but
since validation and generation are offline batch operations, this cost is
acceptable.
